\section{Conclusion}
\label{sec:conclusion}

A lossless air--dielectric--air slab was analyzed using both an equivalent
transmission-line model and a CST frequency-domain simulation. The slab thickness was
approximately one free-space wavelength at 1~GHz, and relative permittivity was swept
from 1 to 10 in 0.25 increments.

The analytical model predicts matched conditions at
$\varepsilon_r=(m/2)^2$. The five values within the sweep range are 1, 2.25, 4, 6.25,
and 9. The preserved CST $|S_{11}|$ response contains minima at the same five values,
which is the central validation result. Field plots at $\varepsilon_r=4$ are consistent
with the intended transverse propagation environment and the integer half-wavelength
condition.

The report does not interpret the exact null depths as convergence-qualified values.
Finite transverse geometry, approximate side boundaries, numerical ports, finite mesh,
rounded slab length, and a recorded port-mode warning distinguish the CST model from
the ideal infinite plane-wave problem. Future work should preserve the native model and
raw complex data, establish mesh and port convergence, examine modal purity, include
loss and oblique incidence, and perform experimental reflection measurements.
